\documentclass[11pt,a4paper]{article}

% ---------------------------------------------------------------------------
% Informe técnico — Proyecto Taxalia
% Compilar con: pdflatex informe.tex  (dos pasadas para el índice)
% ---------------------------------------------------------------------------

\usepackage[utf8]{inputenc}
\usepackage[T1]{fontenc}
\usepackage[spanish,es-tabla]{babel}
\usepackage[a4paper,margin=2.4cm]{geometry}
\usepackage{lmodern}
\usepackage{microtype}
\usepackage{xcolor}
\usepackage{booktabs}
\usepackage{longtable}
\usepackage{array}
\usepackage{enumitem}
\usepackage{titlesec}
\usepackage{fancyhdr}
\usepackage{listings}
\usepackage[hidelinks]{hyperref}

% --- Paleta (dorado/azul corporativo Taxalia) -----------------------------
\definecolor{tax-gold}{HTML}{B8860B}
\definecolor{tax-ink}{HTML}{1F2A37}
\definecolor{tax-slate}{HTML}{475569}
\definecolor{tax-bg}{HTML}{F8FAFC}
\definecolor{tax-line}{HTML}{E2E8F0}

\hypersetup{
  pdftitle={Informe técnico — Proyecto Taxalia},
  pdfauthor={Equipo Taxalia},
  colorlinks=true,
  linkcolor=tax-ink,
  urlcolor=tax-gold,
  citecolor=tax-gold
}

% --- Títulos --------------------------------------------------------------
\titleformat{\section}
  {\color{tax-ink}\Large\bfseries}{\color{tax-gold}\thesection}{0.8em}{}
\titleformat{\subsection}
  {\color{tax-ink}\large\bfseries}{\color{tax-gold}\thesubsection}{0.7em}{}
\titlespacing*{\section}{0pt}{2.2ex plus 1ex minus .2ex}{1.2ex plus .2ex}

% --- Cabecera/pie ---------------------------------------------------------
\pagestyle{fancy}
\fancyhf{}
\renewcommand{\headrulewidth}{0.4pt}
\renewcommand{\footrulewidth}{0pt}
\fancyhead[L]{\small\color{tax-slate}Taxalia}
\fancyhead[R]{\small\color{tax-slate}Informe técnico}
\fancyfoot[C]{\small\color{tax-slate}\thepage}

% --- Estilo de código -----------------------------------------------------
\lstdefinestyle{tax}{
  basicstyle=\ttfamily\small,
  backgroundcolor=\color{tax-bg},
  frame=single,
  rulecolor=\color{tax-line},
  keywordstyle=\color{tax-gold}\bfseries,
  commentstyle=\color{tax-slate}\itshape,
  stringstyle=\color{tax-ink},
  breaklines=true,
  showstringspaces=false,
  columns=fullflexible,
  numbers=none,
  xleftmargin=8pt,
  xrightmargin=8pt,
  aboveskip=10pt,
  belowskip=6pt
}
\lstset{style=tax}

% Comando para identificadores inline
\newcommand{\code}[1]{\texttt{\small #1}}

% Caja de destacado
\newcommand{\destacado}[1]{%
  \begin{center}
  \fcolorbox{tax-line}{tax-bg}{%
    \parbox{0.94\linewidth}{\vspace{4pt}\small #1\vspace{4pt}}}
  \end{center}}

\setlist[itemize]{leftmargin=1.4em,itemsep=2pt,topsep=3pt}

% ===========================================================================
\begin{document}

% --- Portada --------------------------------------------------------------
\begin{titlepage}
\thispagestyle{empty}
\centering
\vspace*{2.5cm}
{\color{tax-gold}\rule{\linewidth}{2pt}}\\[0.6cm]
{\Huge\bfseries\color{tax-ink} Proyecto Taxalia}\\[0.4cm]
{\LARGE\color{tax-slate} Web corporativa + asistente conversacional Lexi}\\[0.3cm]
{\color{tax-gold}\rule{\linewidth}{2pt}}\\[1.2cm]

{\large\bfseries Informe técnico}\\[0.3cm]
{\large Decisiones, tecnología y el sistema multiagente}\\[2.5cm]

\begin{tabular}{rl}
  \textbf{Proyecto:}    & Taxalia (firma de asesoría) \\[3pt]
  \textbf{Arquitectura:}& Meta-repo: frontend Astro + backend de chatbot \\[3pt]
  \textbf{Agente:}      & Lexi — orquestador multiagente sobre Ollama \\[3pt]
  \textbf{Fecha:}       & Junio de 2026 \\
\end{tabular}

\vfill
{\small\color{tax-slate} Documento generado a partir del estado real del
repositorio: ADRs, artefactos SDD y código fuente del backend y el frontend.}
\end{titlepage}

% --- Índice ---------------------------------------------------------------
\tableofcontents
\newpage

% ===========================================================================
\section{Resumen ejecutivo}

Taxalia es la web corporativa de una firma de asesoría profesional
(consultoría, valoración y acompañamiento estratégico). El proyecto va más
allá de un sitio estático: incorpora \textbf{Lexi}, un asistente
conversacional propio construido como un \textbf{sistema multiagente} que se
ejecuta contra un modelo de lenguaje local (Ollama) y se despliega también en
la nube.

La pieza diferencial no es ``un chatbot'' sino la \textbf{ingeniería que lo
sostiene}: un orquestador que decide qué agentes responden, ejecución en
paralelo con cancelación cooperativa, un sintetizador que funde varias
respuestas en una sola, políticas de conducta versionadas en disco y una
estrategia de pruebas y observabilidad de nivel producción.

\destacado{\textbf{En una frase:} una web ligera en Astro con coste de
dependencias casi nulo, y detrás un backend de agentes en TypeScript estricto,
diseñado con metodología \emph{Spec-Driven Development} (SDD), decisiones
documentadas como ADR y un contrato HTTP/SSE que desacopla por completo el
front del motor de IA.}

\subsection{Cifras del proyecto}

\begin{center}
\begin{tabular}{@{}lr@{}}
\toprule
\textbf{Indicador} & \textbf{Valor} \\
\midrule
Código fuente del backend (TypeScript) & $\sim$2.655 líneas \\
Ficheros fuente del backend            & 36 módulos \\
Suites de pruebas del backend          & 22 ficheros (unitarias + integración) \\
Código fuente del frontend             & $\sim$3.220 líneas \\
Agentes especializados                 & 3 (advisory, valuation, financial) \\
Políticas de conducta                  & 5 (cargadas en arranque) \\
Capacidades (skills)                   & 3 \\
Nuevas dependencias npm en el frontend & 0 \\
Decisiones de arquitectura formales    & 2 ADR + 6 SDD \\
\bottomrule
\end{tabular}
\end{center}

% ===========================================================================
\section{Contexto y objetivos}

La web de Taxalia se diseñó para presentar la identidad de la empresa, mostrar
sus servicios, publicar contenidos de blog y facilitar el contacto, con
soporte multilingüe (inglés y español). El sitio es estático, ligero y fácil
de desplegar.

El punto de partida tenía una limitación clara: el \emph{widget} de chat era un
ancla estática que redirigía a la página de contacto. El objetivo del trabajo
fue \textbf{sustituir ese marcador por una superficie conversacional real},
sin comprometer la simplicidad del sitio: sin integraciones de Astro, sin
cambios en el \emph{pipeline} de build y \textbf{sin añadir una sola
dependencia npm al frontend}.

\destacado{La restricción de diseño \emph{``cero dependencias nuevas en el
front''} obligó a una solución limpia: el cliente de chat usa solo APIs nativas
del navegador (\code{fetch}, \code{ReadableStream}, \code{AbortController},
\code{crypto.randomUUID}). Toda la complejidad vive en un backend
independiente, conectado por un contrato HTTP.}

% ===========================================================================
\section{Arquitectura general}

El repositorio es un \textbf{meta-repo} con dos proyectos hermanos y autónomos,
más la tooling transversal:

\begin{lstlisting}
taxalia/                  <- meta-repo (sin package.json propio)
|-- docs/adr/             <- decisiones de arquitectura (ADR)
|-- openspec/             <- artefactos SDD (propuesta, diseno, specs, tareas)
|-- scripts/setup.mjs     <- bootstrap de ambos lados (cross-platform)
|-- setup.sh              <- wrapper Unix delgado
|-- frontend/             <- sitio Astro (submodulo Git)
|   |-- src/components/    ChatWidget.astro, Header, Hero, Services...
|   |-- src/scripts/       chat-client.ts (consumidor SSE)
|   `-- src/pages/         index, services, blog, contact, /es/...
`-- backend/              <- chatbot multiagente (submodulo Git)
    `-- src/{chat,agents,skills,conducta,dispatch,ollama,observability}
\end{lstlisting}

Cada lado es independiente: \code{cd frontend \&\& npm install \&\& npm run
dev} funciona sin asumir nada de la raíz, y el backend declara su propia
configuración de TypeScript sin depender de ficheros del front.

\subsection{Arquitectura ``screaming''}

El backend aplica \emph{screaming architecture}: los nombres de carpeta gritan
el dominio, no el framework. Las carpetas (\code{chat}, \code{agents},
\code{skills}, \code{conducta}, \code{dispatch}, \code{ollama},
\code{observability}) sobreviven a cualquier cambio de tecnología: si se
sustituye Hono, \code{chat/} permanece; si se sustituye Ollama, \code{ollama/}
pasa a llamarse \code{inference/}. El dominio manda sobre la herramienta.

% ===========================================================================
\section{Decisiones tomadas}

Las decisiones no son implícitas: están documentadas como \textbf{ADR}
(Architecture Decision Records) y como resoluciones dentro del flujo
\textbf{SDD}. Esto deja una traza auditable del \emph{por qué}, no solo del
\emph{qué}.

\subsection{Decisiones de estructura (ADR)}

\begin{center}
\renewcommand{\arraystretch}{1.3}
\begin{tabular}{@{}p{0.16\linewidth}p{0.46\linewidth}p{0.30\linewidth}@{}}
\toprule
\textbf{ADR} & \textbf{Decisión} & \textbf{Razón} \\
\midrule
0001 — Backend hermano &
Mantener backend como carpeta hermana \code{backend/}, sin npm workspaces. &
Un solo servicio; se difiere la tooling de monorepo hasta que aparezca un
segundo backend o un paquete compartido. \\
0002 — Frontend agrupado &
Mover el proyecto Astro a \code{frontend/}; la raíz queda como meta-repo puro
sin código de aplicación. &
Separa la app de la tooling transversal; rompe el acoplamiento oculto del
\code{tsconfig} del backend hacia el front. \\
\bottomrule
\end{tabular}
\end{center}

Ambos ADR fijan el mismo \textbf{disparador de migración} a npm workspaces:
cuando exista un segundo backend \emph{o} un paquete compartido por ambos
lados. Hasta entonces, la forma multirepo es la correcta.

\subsection{Decisiones del agente (SDD)}

El sistema multiagente se diseñó con \emph{Spec-Driven Development}: primero
exploración, luego propuesta, especificaciones, diseño y tareas, antes de
escribir código. Las decisiones clave resueltas:

\begin{center}
\renewcommand{\arraystretch}{1.3}
\begin{tabular}{@{}p{0.22\linewidth}p{0.40\linewidth}p{0.30\linewidth}@{}}
\toprule
\textbf{Cuestión} & \textbf{Decisión} & \textbf{Razón} \\
\midrule
Agentes v1 &
\code{advisory}, \code{valuation}, \code{financial}. &
Alinea los agentes con la taxonomía real de servicios de la web. \\
Conducta v1 &
5 políticas: no fingir, citar fuentes, respuesta bilingüe, sin PII, traspaso a
humano. &
Cubre la superficie de riesgo v1: alucinación, citación, idioma, privacidad y
escalado. \\
Concurrencia &
Tope de 2 \emph{dispatches} simultáneos vía semáforo en proceso. &
Un equipo M3 de 16\,GB se satura con 2 dispatches; más invita a OOM. \\
Reintentos &
Sin reintentos automáticos ante fallo transitorio; \emph{fail loud} con aviso. &
Los reintentos silenciosos en un modelo local pequeño producen UX confusa. \\
Autenticación &
Anónimo + lista blanca CORS. &
El backend corre en local sobre Ollama; sin exposición pública en v1. \\
Estimador de tokens &
Heurística calibrada (4 chars/token + sobrecarga). &
Validada contra \code{prompt\_eval\_count} real de Ollama, error $\leq$15\%. \\
\bottomrule
\end{tabular}
\end{center}

\subsection{Decisión de entrega: PRs encadenados}

El trabajo se planificó como \textbf{5 pull requests encadenados}, cada uno
con un presupuesto de revisión $\leq$400 líneas, para mantener la revisión
enfocada: (1) higiene + setup, (2) esqueleto Hono + contrato SSE,
(3) cargadores + ensamblado del system prompt, (4) orquestador + dispatch
paralelo + sintetizador, (5) reemplazo del widget y cableado SSE en el front.
Cada PR define su propio plan de \emph{rollback}.

% ===========================================================================
\section{Stack tecnológico}

\subsection{Frontend}

\begin{center}
\renewcommand{\arraystretch}{1.25}
\begin{tabular}{@{}p{0.26\linewidth}p{0.66\linewidth}@{}}
\toprule
\textbf{Tecnología} & \textbf{Rol} \\
\midrule
Astro 6.3 & Generación de sitio estático, ligero y fácil de desplegar. \\
TypeScript & \code{chat-client.ts}: consumidor SSE y enlaces al DOM. \\
i18n manual & Estrategia propia (\code{/es/} + \code{i18n.ts}), sin
\code{@astrojs/i18n}. \\
CSS propio & \code{lb-co.css}; estética corporativa dorado/azul. \\
APIs nativas & \code{fetch}, \code{ReadableStream}, \code{AbortController}. \\
\bottomrule
\end{tabular}
\end{center}

\subsection{Backend del agente}

\begin{center}
\renewcommand{\arraystretch}{1.25}
\begin{tabular}{@{}p{0.26\linewidth}p{0.66\linewidth}@{}}
\toprule
\textbf{Tecnología} & \textbf{Rol} \\
\midrule
Node 24 + TypeScript estricto & Runtime y tipado fuerte de extremo a extremo. \\
Hono & Framework HTTP que sirve la superficie de chat y el streaming SSE. \\
Zod & Esquemas como única fuente de verdad (request, SSE, defs de ficheros). \\
Ollama (\code{ollama-js}) & Cliente del modelo; \code{gemma4:e4b} en local,
\code{gemma4:31b-cloud} en producción. \\
Pino & Logging estructurado con \code{requestId} por petición. \\
Vitest & Pruebas unitarias y de integración (proyectos separados). \\
\bottomrule
\end{tabular}
\end{center}

El \textbf{contrato de tipos} se define una sola vez con Zod y se infiere con
\code{z.infer<...>}, de modo que el esquema de validación en runtime y el tipo
de TypeScript nunca divergen.

% ===========================================================================
\section{El agente Lexi: el corazón del proyecto}

Aquí está el verdadero valor. ``Lexi'' no es una llamada a un LLM con un
prompt: es un \textbf{sistema multiagente} con orquestación, paralelismo,
síntesis, políticas y observabilidad. Estos son sus componentes.

\subsection{Personas como ficheros en disco}

Agentes, capacidades y conducta son \textbf{ficheros Markdown con frontmatter
YAML}. Cambiar el comportamiento del asistente no requiere tocar código:
basta editar un fichero y reiniciar (o llamar a \code{POST /admin/reload} en
desarrollo).

\begin{lstlisting}
---
id: advisory
name: Advisory Services Assistant
description: Helps users understand Taxalia advisory services.
system_prompt: |
  You are Taxalia's advisory services assistant. Explain
  engagement models clearly in English or Spanish...
tools: [lookup-engagement-model, capture-lead]
tags: [advisory, services]
---
\end{lstlisting}

Esta separación \emph{datos vs. código} es lo que convierte al asistente en un
producto mantenible por personas no programadoras.

\subsection{El pipeline de una conversación}

Cada mensaje del usuario atraviesa cuatro etapas claramente separadas:

\begin{enumerate}[label=\textbf{\arabic*.}]
\item \textbf{Orquestación} — una llamada al modelo en formato JSON decide
\emph{qué} agentes deben responder (\code{agentsToRun}). Saludos o intención
desconocida producen un array vacío. Nunca se inventan agentes.
\item \textbf{Dispatch paralelo} — los agentes seleccionados se ejecutan a la
vez (\code{Promise.allSettled}), cada uno con timeout de 30\,s y su propio
\code{AbortController} enlazado a la señal externa.
\item \textbf{Síntesis} — si responden 2 o más agentes, un sintetizador funde
sus salidas en una única respuesta coherente.
\item \textbf{Streaming SSE} — la respuesta vuelve al navegador token a token
como \emph{Server-Sent Events}; el widget va componiendo la burbuja.
\end{enumerate}

\destacado{\textbf{Cancelación cooperativa de extremo a extremo.} Si el usuario
cierra el chat, el \code{AbortController} del navegador propaga la señal a la
ruta de Hono, de ahí a cada \code{AbortController} por agente y de ahí al
\emph{stream} de Ollama. Todo se cancela en $\leq$200\,ms. Es el tipo de detalle
que separa un prototipo de un sistema bien construido.}

\subsection{Ensamblado del system prompt por capas}

El prompt de sistema se construye en un orden \textbf{deliberado}, de lo más
general a lo más específico: conducta $\rightarrow$ persona del agente
$\rightarrow$ metadatos de skills $\rightarrow$ historial. Este orden no es
estético: es una defensa empírica frente a \emph{prompt injection}. Además,
si el prompt supera 1.500 tokens, se rechaza con error explícito en lugar de
degradar en silencio.

\subsection{Las 5 políticas de conducta}

Las políticas se cargan en arranque y se concatenan en \emph{cada} prompt de
agente. El arranque \textbf{verifica que sean exactamente 5} (una sexta haría
fallar el boot, a propósito: cada política debe ser una incorporación
consciente).

\begin{center}
\renewcommand{\arraystretch}{1.25}
\begin{tabular}{@{}p{0.26\linewidth}p{0.62\linewidth}@{}}
\toprule
\textbf{Política} & \textbf{Regla} \\
\midrule
\code{never-pretend}      & Decir ``no lo sé'' antes que inventar hechos,
                            precios o conclusiones legales. \\
\code{cite-sources}       & Citar una página de servicio de Taxalia o admitir
                            que no hay fuente. \\
\code{bilingual-response} & Responder en el idioma del usuario; no mezclar. \\
\code{privacy-no-pii}     & No almacenar, repetir ni inventar datos personales. \\
\code{handoff-to-human}   & Si piden un humano, ofrecer el enlace de contacto. \\
\bottomrule
\end{tabular}
\end{center}

\subsection{Observabilidad y manejo de fallos}

El sistema es transparente sobre lo que ocurre:

\begin{itemize}
\item \textbf{Trazabilidad:} cada petición lleva un \code{requestId} (cabecera
\code{X-Request-Id} o UUID v4 generado) que viaja por todos los logs y vuelve
en el evento final del SSE.
\item \textbf{Métricas:} contadores en memoria expuestos en \code{GET /metrics}
(chats iniciados/completados, timeouts por agente, fallos parciales, errores de
parseo del orquestador\dots).
\item \textbf{Fallo parcial con aviso:} si un agente falla, su error se reporta
en \code{done.agents[]} y el sintetizador trabaja con los que sí respondieron;
el usuario ve un aviso de ``respuesta posiblemente incompleta''. Nunca se oculta
un fallo.
\item \textbf{Privacidad por diseño:} los logs registran \code{requestId},
agentes, idioma, conteo de tokens y latencia, \textbf{nunca} el contenido de
los mensajes.
\end{itemize}

% ===========================================================================
\section{Calidad: pruebas y presupuestos}

\subsection{Estrategia de pruebas}

22 suites cubren las dos capas. Lo destacable es \emph{qué} se prueba:

\begin{itemize}
\item \textbf{Pureza del system prompt:} se ensambla 1.000 veces y se verifica
que el resultado es idéntico byte a byte.
\item \textbf{Carrera de cancelación:} un cliente sintético abre la conexión,
espera el primer evento y aborta; se afirma que todos los \code{AbortController}
disparan en $\leq$200\,ms.
\item \textbf{Tasa de parseo del orquestador:} el modelo real se enfrenta a 20
mensajes de prueba; criterio de aprobado $\geq$16/20 con JSON parseable.
\item \textbf{Arranque en frío:} la primera petición contra un Ollama recién
levantado responde dentro de la ventana de gracia de 60\,s.
\end{itemize}

\subsection{Presupuestos de rendimiento}

Las metas no son aspiracionales: están escritas y testeadas.

\begin{center}
\renewcommand{\arraystretch}{1.25}
\begin{tabular}{@{}p{0.58\linewidth}r@{}}
\toprule
\textbf{Métrica} & \textbf{Presupuesto} \\
\midrule
\code{GET /health} (sin llamar a Ollama)            & $\leq$100\,ms \\
Gracia de arranque en frío (una vez por proceso)    & 60\,s \\
Timeout por agente                                  & 30\,s \\
System prompt ensamblado                            & $\leq$1.500 tokens \\
Propagación del abort (cliente $\rightarrow$ Ollama)& $\leq$200\,ms \\
Primer token con 2 agentes (en caliente)            & $\leq$8\,s \\
Respuesta completa, 2 agentes ($\leq$200 tokens)    & $\leq$30\,s \\
\bottomrule
\end{tabular}
\end{center}

% ===========================================================================
\section{Despliegue}

El sistema funciona en dos modos sin cambios de código, solo de configuración:

\begin{itemize}
\item \textbf{Local:} Ollama en \code{127.0.0.1:11434} con el modelo
\code{gemma4:e4b}. El script \code{setup.sh} / \code{scripts/setup.mjs} es
idempotente: comprueba Ollama, descarga el modelo y ejecuta \code{npm ci} en
ambos lados.
\item \textbf{Producción:} backend como función \emph{serverless} en Vercel
(\code{vercel.json} incluye los ficheros de agentes/skills/conducta en el
bundle), apuntando al modelo en la nube \code{gemma4:31b-cloud}.
\end{itemize}

La selección de host y modelo según \code{NODE\_ENV} está centralizada en
\code{config.ts} y validada con Zod, de modo que un entorno mal configurado
falla en el arranque con un mensaje claro, no a mitad de una petición.

% ===========================================================================
\section{Estado actual y próximos pasos}

El núcleo conversacional está operativo de extremo a extremo: orquestación,
dispatch paralelo, síntesis, streaming, conducta, observabilidad, pruebas y
despliegue dual. El frontend ya integra el widget real con avatar, redimensionado
y página de \emph{health check}.

Quedan tareas de producto recogidas en el \emph{backlog} del proyecto, entre
ellas: completar páginas internas (servicios individuales, blog con
\emph{Content Collections}), conectar el formulario de contacto a un backend
real, cerrar la internacionalización al español, y reforzar SEO/accesibilidad
(sitemap, \code{robots.txt}, Open Graph, JSON-LD, navegación por teclado).

% ===========================================================================
\section{Conclusión}

Taxalia demuestra que se puede añadir IA conversacional a un sitio estático
\textbf{sin sacrificar su simplicidad ni inflar sus dependencias}. El valor no
está en ``tener un chatbot'', sino en cómo está construido:

\begin{itemize}
\item Un \textbf{sistema multiagente} real —orquestador, paralelismo,
sintetizador— en lugar de una llamada suelta a un LLM.
\item \textbf{Separación datos/código:} personas y políticas como ficheros
editables sin tocar el motor.
\item \textbf{Robustez de producción:} cancelación cooperativa, fallo parcial
con aviso, presupuestos de rendimiento testeados y observabilidad con
\code{requestId} de punta a punta.
\item \textbf{Disciplina de ingeniería:} decisiones documentadas (ADR + SDD),
TypeScript estricto, contrato de tipos único con Zod y entrega en PRs
encadenados revisables.
\end{itemize}

\destacado{El agente montado es la mejor carta de presentación técnica del
proyecto: un componente pequeño en líneas de código pero grande en criterio de
ingeniería, donde cada decisión está justificada y cada modo de fallo está
previsto.}

\end{document}
